\chapter{Davi e Golias: passado, presente e futuro do Vale do Silício}

\epigraph{
\justifying{
	Eles vão.
	Eles abandonam Omelas, caminham rumo à escuridão, e nunca voltam.
	O lugar para onde vão é ainda menos imaginável para a maioria de nós do que a cidade da felicidade.
	Eu não consigo nem descrevê-lo.
	É possível que não exista.
	Mas eles parecem saber aonde vão, aqueles que fogem de Omelas.
}	
}{--- Ursula K. Le Guin, \textit{``Aqueles que fogem de Omelas''}}

No curso desta monografia, as figuras do algoritmo e dos dados foram tratadas não como meras ideias abstratas, mas sim como ferramentas materiais que têm consequências tangíveis.
Tratá-los individualmente, como entidades isoladas, é insuficiente para compreender sua real importância; é necessário situá-los dentre um conjunto de estruturas sociais, analisando sua função social, econômica e política.
Esse terceiro capítulo colocará uma dimensão histórica nessa análise, tentando assim entender \textit{como} que a nossa sociedade se organizou desse jeito.

O desenvolvimento tecnológico está sempre inserido em um universo de atores sociais, ações, e mecanismos, que o motivam e possibilitam.
A Segunda Revolução Industrial, no século XIX, aconteceu mediante uma instauração do capitalismo como sistema hegemônico e consolidação da burguesia como classe dominante; a assim chamada ``Revolução Verde'', isto é, a substituição das práticas campesinas tradicionais por uma agricultura mecanizada e moderna, levou a um amplo êxodo rural, pobreza no campo e concentração de terras [citação necessária]; a Guerra Fria do século XX, um conflito de natureza geopolítica, gerou enormes avanços tecnológicos, como a bomba atômica.
Vistos sozinhos, soltos e ao léu, tanto o surgimento das indústrias, quanto a mecanização da agricultura, quanto o estudo atômico são impressionantes e deveras importantes para a humanidade.
Porém, eles existem dentro de estruturas sociais, e são, por isso, malfeitores para alguns, e coadjuvantes para outros, seja em uma disputa por poder, terras, dinheiro, entre outros.

E isso não é diferente com os algoritmos e os dados.
Retome a ``retórica do Vale do Silício'' mencionada no primeiro capítulo: é, por vezes, um discurso de emancipação social por meio das plataformas, mas que ignora as condições estruturais das relações de trabalho e poder.
A própria narrativa fatalista de que as máquinas vão dominar a humanidade, \textit{a lá} Matrix, cai nessa falácia.
Ela entende as máquinas não como produtos de uma sociedade, criadas por pessoas com intenções e motivações próprias, mas sim como interventoras quase divinas.
E isto não é verdade: toda e qualquer máquina está integrada na sociedade que a produz.
Mitos como o da IA que recusa a desligar-se, são falsos de um ponto de vista técnico \textit{e social}.
Todas essas novas tecnologias são ferramentas de uma sociedade, de seu sistema de produção, e de sua política; a ameaça que elas representam não é algo inexorável, da sua natureza, mas sim, uma ameaça que converge com os interesses de um grupo de atores sociais humanos, feitos de carne e osso.

Por isso, esse capítulo tratará do assunto em três tempos, começando com uma análise da história da Internet, de modo a entender como as estruturas sociais de vigilância consolidaram-se ao redor dela.
Em seguida, serão analisadas as tentativas de regulação do meio \textit{online}, onde o poder público, antes ausente desse espaço, passa a tentar exercer poder e soberania sobre ele.
Por fim, será apresentado um modo de organização diferente, que destoa dos monopólios e da imposição assimétrica da vigilância que tanto se falou sessa monografia: o \textit{software livre}.

\section{A terra prometida}

A \textit{world wide web} (WWW), que, ao pé da letra, significa ``rede de alcance global'' é uma suíte de protocolos e normas, construída sobre uma infraestrutura de comunicação entre computadores que hoje chamamos de Internet\footnote{A ideia e implementação de redes de computadores interligados precede a WWW, já existindo, por exemplo, na forma da ARPANET americana e da CYCLADES francesa}.
% Ela consiste em quatro componentes principais: o HTML (\textit{HyperText Markup Language}, ou Linguagem de Marcação de HiperTexto), um formato de arquivo pelo qual documentos deveriam ser compartilhados, o HTTP (\textit{HyperText Transfer Protocol}, ou Protocolo de Transferência de Hipertexto), meio pelo qual esses documentos são compartilhados, os URIs (\textit{Uniform Resource Identifiers}, ou Indicadores Uniformes de Recurso), pelos quais os documentos seriam identificados, e os navegadores, pelos quais os documentos seriam visualizados \cite{niaz_webbased_2022}.
Ela foi criada pelo britânico Tim Berners-Lee em 1989, para facilitar a troca de informações entre laboratórios e computadores da Organização Europeia para a Investigação Nuclear, conhecida também por sua sigla em francês CERN, e rapidamente se espalhou por outras instituições científicas.
Os programas e protocolos que constituíam essa rede foram publicados abertamente, e, em pouco tempo, a WWW se espalhou para além do meio acadêmico.

De início, a WWW (que, apesar de não serem iguais, será a partir de agora referida também como Internet, para facilitar a notação) era verdadeiramente um espaço livre.
O indivíduo, se pertencesse ao pequeno grupo daqueles que tinham acesso a ela, tinha total autonomia sobre seu uso, sem responder à nenhuma outra regulação senão a das comunidades em que se inseria.
A ``Declaração da Independência do Ciberespaço'' \cite{barlow_declaration_1996}, apresentada por John Perry Barlow ao Fórum Econômico Mundial em Davos, Suíça, no ano de 1996, é bastante sintética em relação a essa autoproclamada ``independência'' da Internet.
% \footnote{Original: \textit{``
% Governments of the Industrial World, you weary giants of flesh and steel, I come from Cyberspace, the new home of Mind. On behalf of the future, I ask you of the past to leave us alone. You are not welcome among us. You have no sovereignty where we gather.
% We have no elected government, nor are we likely to have one, so I address you with no greater authority than that with which liberty itself always speaks. I declare the global social space we are building to be naturally independent of the tyrannies you seek to impose on us. You have no moral right to rule us nor do you possess any methods of enforcement we have true reason to fear.
% [\dots]
% We are creating a world that all may enter without privilege or prejudice accorded by race, economic power, military force, or station of birth.
% We are creating a world where anyone, anywhere may express his or her beliefs, no matter how singular, without fear of being coerced into silence or conformity.
% Your legal concepts of property, expression, identity, movement, and context do not apply to us. They are all based on matter, and there is no matter here.
% Our identities have no bodies, so, unlike you, we cannot obtain order by physical coercion. We believe that from ethics, enlightened self-interest, and the commonweal, our governance will emerge.
% [\dots]
% We will create a civilization of the Mind in Cyberspace. May it be more humane and fair than the world your governments have made before.''}}

\begin{quoting}
Governos do Mundo Industrial, gigantes cansados de carne e aço, eu venho do Ciberespaço, o novo lar da Mente.
Em nome do futuro, eu peço que vocês, do passado, nos deixem em paz.
Vocês não são bem vindos entre nós. Vocês não tem soberania onde nos reunimos.

Nós não temos governo eleito, e nem é provável que o tenhamos, então eu me dirijo a vocês sem autoridade maior do que aquela com que a liberdade sempre fala.
Eu declaro que o espaço social global que estamos construindo é naturalmente independente das tiranias que vocês buscam nos impor.
Vocês não tem o direito moral de nos governar, nem possuem quaisquer métodos de imposição que tenhamos razão verdadeira para temer.

[...]

Nós estamos criando um mundo em que todos poderão entrar sem privilégio ou preconceito de raça, poder econômico, força militar, ou origem.

Nós estamos criando um mundo onde qualquer pessoa, em qualquer lugar, pode expressar suas crenças, não importa quão singulares, sem medo de serem coagidas ao silêncio ou à conformidade.

Seus conceitos legais de propriedade, expressão, identidade, movimento, e contexto, não se aplicam a nós.
Eles são todos baseados em matéria, e não há matéria aqui.

Nossas identidades não têm corpos, por isso, diferente de vocês, nós não podemos obter ordem por coerção física.
Nós acreditamos que pela ética, pelo interesse próprio iluminado, e pelo bem-comum, nossa governança vai emergir.

[...]

Nós estamos criando uma civilização da Mente no Ciberespaço.
Que ela seja mais humana e justa do que o mundo que seus governos produziram.

\cite[tradução~própria]{barlow_declaration_1996}
\end{quoting}

Nessa concepção, o ``ciberespaço'', como Barlow o chama, seria completamente segregado e independente do poder público e do Estado.
Ele seria isonômico, ignorante da realidade física, e atuaria em contraposição aos governos pré-estabelescidos.
Como ``o novo lar da Mente'', ele seria a mais nova etapa do desenvolvimento da civilização; existe nesse trecho uma noção evolucionária, que levaria a um espaço (o digital) em que os indivíduos pudessem interagir anonimamente com quase absoluta liberdade.
% Não é difícil conectar o pensamento de Barlow ao Iluminismo de Kant: ambos tratam de um espaço público onde haveria troca irrestrita de opiniões entre indivíduos autônomos, e através disso, uma emancipação da sociedade em geral.
O ciberespaço de Barlow teria a anonimidade como imposição mesológica, e se encontraria completamente separado (e não integrado) do Estado, negando o que poderia chamar-se de política tradicional, e substituindo-a com a nova governança do ciberespaço.
Alguns dos teóricos da época previam uma ágora digital, isto é, uma democracia direta fundada em meios digitais.

Esse trecho demonstra o que Paris Marx chama de \textit{ciberlibertariansmo}: a principal doutrina a guiar a Internet em sua concepção.
Não é coincidência que esse documento, delimitador de uma suposta sociedade alternativa que afronta a ordem vigente, foi apresentado ao mais alto escalão dos governos e do empresariado, no Fórum Econômico Mundial.
Ele foi, por muito tempo, o discurso oficial do mundo digital \cite{marx_pavel_2025}.

O ciberlibertarianismo é, desde sua concepção, uma doutrina decididamente idealista.
Ela idealiza o espaço digital, que aparece, totalmente democrático e livre, em uma relação mutualmente ininfluenciável com o mundo real.
Tão democrático que rejeita que quaisquer variáveis do mundo material -- dinheiro, por exemplo -- possam influenciá-lo.

Há, para \textcite{marx_pavel_2025} uma continuidade entre a retórica ciberlibertária e o discurso ideológico do Vale do Silício, assim como foi discutido no primeiro capítulo com base no texto de Morozov.
O autor analisa, em especial, a relação entre essa retórica e o conceito de ``liberdade de expressão'', hoje apropriado, de modo distorcido, pela extrema-direita.
O espaço digital de Barlow, afinal, é totalmente livre, nesse sentido.
Podem-se expressar quaisquer ideias, ``não importa o quão singulares''.
É uma retórica de emancipação, que busca criar um mundo separado, e nisso, ignora as conexões do digital com o físico.
Discursos de ódio, neonazismo, e pornografia infantil são, nessa concepção, viabilizados \cite{morozov_big_2018}.

Apesar de ter uma dimensão em algum grau louvável de democratização do debate público, a doutrina de Barlow é falha.
E, além disso, possibilitou a criação de uma Internet que, hoje, cerceia, vigia, e influencia, ora autoritariamente, seus usuários.

\section{Gigantes de aço}
Todo o ódio dos ciberlibertários contra o poder estabelecido era dirigido ao Estado, e fim da história.
Não se disse um pio sobre o poder privado, que prontamente viu uma oportunidade de crescer seu mercado, e se proliferou livremente até ocupar quase toda a Internet.
Formaram-se grandes monopólios privados, que, devido ao caráter oficial da separação entre o material e o digital, permaneceram por muito tempo completamente desreguladas.
Para \textcite{nunes_coase_2025},

\begin{quoting}
Em paralelo ao surgimento de milhões de microempresas de baixo capex [despesas de capital], essas novas tecnologias permitiram a eclosão de uma dezena de corporações colossais, cujas dimensões e funções paraestatais só encontram paralelo nas companhias das índias holandesas da expansão ultramarina. A Nvidia, com US\$ 4,4 trilhões de valor de mercado, se equipara ao PIB da Alemanha. A Microsoft e Apple, ambas com capitalização de mercado acima de US\$ 3 trilhões, aproximam-se do PIB do Reino Unido. Amazon e Google, com cerca de US\$ 2,5 trilhões, superam economias como a italiana e a canadense. A Meta, em torno de US\$ 1,9 trilhão, equivale ao Brasil, enquanto a Tesla, com US\$ 1,1 trilhão, se aproxima de países de porte médio.
\end{quoting}

Hoje, a gigantesca maioria do tráfego na Internet é realizado nos domínios de poucas empresas, que também controlam a infraestrutura que a sustenta.
A dissociação do mundo digital beneficiou diretamente essas empresas, ao permitir que elas tomassem para si as terras digitais.
Essa doutrina tinha um envase anti-sistema, mas seguia à risca o princípio neoliberal de não intervenção do Estado na economia e desregulamentação do mercado.
As empresas do ramo, por isso, gozaram, durante muito tempo, de quase completa liberdade econômica, pouca carga tributária e uma total dominação de seus mercados \cite[p.~156]{morozov_big_2018}, transformando-se assim nessas corporações colossais.

Criaram-se gigantes inabaláveis, sentados sobre tronos feitos de dados pessoais coletados ao longo de múltiplos anos, que, por meio do aprendizado de máquina, puderam expressar-se na forma de modelos computacionais dos indivíduos vigiados.
O poder sobre os dados igualou-se, como já foi dito na seção \ref{sec:1:politica}, a poder político; não é à toa que \textcite{nunes_coase_2025} os equaciona, em termos de capital, a países inteiros.
E esse poder político não se manifesta somente em interferência política e lobbying, mas também, por exemplo, ao influenciar a direção das pesquisas científicas, direcionando-as segundo seus interesses.
De acordo com \textcite{kalluriComputervisionResearchPowers2025}, de 1990 até 2010 a quantidade de artigos ligados a patentes relacionadas a vigilância por câmeras de seres humanos quintuplicou.
Os autores também apontam que há um ofuscamento linguístico nesses artigos, fazendo-os utilizar palavras menos carregadas -- trocando, por exemplo, ``humano'' por ``objeto'' -- para abstrair a realidade que eles coadjuvam. 
O projeto político de uma sociedade de vigilância ganha, deste modo, capacidade tecnológica para se concretizar. 

Contraditoriamente, o mercado digital \textit{laissez faire} instituído pelo Vale do Silício dependeu, e continua dependendo, de amplo investimento estatal para existir, seja pela construção de infraestrutura física para a expansão do meio (cabos, torres), isenções fiscais (como vem acontecendo no Brasil com a construção de \textit{Data Centers}), ou até por serem os Estados os principais provedores de pesquisa científica, essencial para a manutenção dessa indústria.
As contribuições estatais são, porém, veementemente ignoradas pelos proponentes dessa ideologia.
O Vale do Silício prefere, nesse sentido, enaltecer a autonomia individual dos empreendedores, que, ignorando completamente a natureza coletiva e cumulativa do conhecimento científico, seriam os verdadeiros criadores da inovação \cite{barbrook_ideologia_2018}.

Também é contraditória a relação dessa doutrina -- a qual \textcite{barbrook_ideologia_2018} chamam de Ideologia Californiana -- com o trabalho.
Ela defende o que os autores chama de ``determinismo tecnológico'': uma noção de que ter sua força de trabalho automatizada é uma previsão fadada a acontecer.
Em outras palavras, é uma doutrina fatalista.
Aqui, a etimologia da palavra ``robô'', cunhada pelo tcheco Karel Čapek com o significado de ``trabalho forçado'', cai bem: deseja-se uma horda de escravos sintéticos, serventes capazes de manter o padrão de vida das elites tecnológicas, pois trabalhadores robóticos não podem nem fazer greve, nem reclamar por salários.
O Vale do Silício almeja uma cristalização da dominação de classes na forma do trabalho automatizado, mas, até lá, não pode prescindir do trabalho humano que mantém essa indústria funcionando.
E nem o poderá no futuro, como argumentam os autores, pois sempre será necessária a criação, modernização e ampliação do maquinário, um trabalho exclusivamente realizado por humanos.

A emancipação da humanidade, assim, não viria pela melhora da qualidade de vida daqueles que hoje sofrem, mas sim por meio de sua eliminação enquanto classe.
A ideologia californiana ignora qualquer tipo de desigualdade hoje existente, embora se baseando nelas, e prevê um futuro melhor, mas excludente; a ágora digital não abre espaço para as populações que, pela lógica do capital, ainda não puderam acessá-la.
O projeto econômico do Vale do Silício não é de todo diferente do neoliberalismo.

É útil aqui tomar uma definição mais restrita de ideologia como um discurso que promove e normaliza uma ordem estabelecida.
A ideologia californiana abarca dentro de si uma oposição dialética entre pregar o ódio aos Estados, à burocracia, e, ao mesmo tempo, fortalecer estruturas de poder nesse mesmo mundo que criticam.
Ela é um excelente exemplo do que Mark Fisher entende como ``realismo capitalista'': a ideia de que o capitalismo assimilou para dentro de si discursos de oposição ao próprio capitalismo.
Os ciberlibertários são capazes de perceber algo de errado com a realidade -- por isso, \textit{realistas} -- mas, presos pelo discurso ideológico, são incapazes de conceber uma alternativa que não seja benéfica para o capitalismo. \cite{fisher_realismo_2020}

Para o autor, os circuitos de entretenimento que se consolidaram pelos canais digitais são, também, uma forma de dominação que, diferentemente das formas anteriores de dominação que aconteciam com coação física no espaço fechado de uma fábrica, agora transformam os indivíduos em seres de consumo completamente tomados por uma \textit{impotência reflexiva}.
Isto é, tem consciência de sua situação precária, mas, concomitantemente, tem consciência de que não há possibilidade de escapar disso.
Foram dominados por uma \textit{pós-lexia}: não é mais necessária a leitura e a interpretação de texto, uma vez que o espaço digital tal qual ele foi consolidado pode ter seus conteúdos absorvidos sem que esses canais sejam empregados \cite{fisher_realismo_2020}.
Tudo movido a dados, claro.

Tornou-se necessário, em algum momento, regulamentar e regularizar o meio digital, principalmente a partir do momento em que ele começou a influenciar o fazer político dos Estados, e desafiou vários pressupostos sobre os quais a lei foi construída, como propriedade privada (por meio da pirataria) e identificabilidade dos cidadãos (por meio do que, \textit{a priori}, era uma anonimização).
Junto, surgiu uma série de problemas provenientes da livre circulação de informações, como discurso de ódio, fraude, comunicação não-autorizada de dados pessoais sensíveis, entre outros.
Isto é, informações que antes passavam pela censura e crivo estatal e então estariam escapando dessa supervisão.

Para muitos desses problemas, os Estados encontraram soluções, vindas principalmente de uma queda no fator da anonimidade.
De acordo com \textcite{han_infocracia_2022}, formou-se uma sociedade transparente, onde os indivíduos voluntariamente dão suas identidades para os serviços que usam.
Com isso, é possível refletir uma ação criminosa digital em uma consequência física para o ator envolvido;  as mesmas abordagens coercivas estatais usadas em meios físicos podem também ser utilizadas em meios digitais.
Quando o objeto da coerção é uma empresa, essa quebra de anonimidade é desnecessária, uma vez que elas sempre foram organizações não-anônimas e com equivalentes em ambos os meios.

No Brasil, há duas leis principais que regulamentam a atuação do Estado frente à Internet.
São elas o Marco Civil da Internet (MCI), de 2014, e a Lei Geral de Proteção de Dados Pessoais (LGPD), de 2018.
O Marco Civil da Internet é, primordialmente, uma lei que dita os princípios que o uso da Internet deverá seguir.
Dentre eles figuram a liberdade de expressão, conforme delimitada pela Constituição Federal, a proteção da privacidade e dos dados pessoais, e a neutralidade de rede\footnote{A neutralidade de rede, como a própria lei define, é a imposição que o ``responsável pela transmissão, comutação ou roteamento tem o dever de tratar de forma isonômica quaisquer pacotes de dados, sem distinção por conteúdo, origem e destino, serviço, terminal ou aplicação''. Isso garante que a infraestrutura que sustenta a Internet é isonômica, e que qualquer tipo de segregação de tráfego ou censura deverá acontecer em um nível maior de abstração.}.
Ele também descreve uma série de direitos dos usuários, o principal deles sendo o direito ao acesso à Internet, e deveres dos provedores.

Para o fim dessa análise, considere os seguintes trechos da lei:

% \begin{quoting}
% Art. 3º A disciplina do uso da internet no Brasil tem os seguintes princípios:

% I - garantia da liberdade de expressão, comunicação e manifestação de pensamento, nos termos da Constituição Federal;

% II - proteção da privacidade;

% III - proteção dos dados pessoais, na forma da lei;

% IV - preservação e garantia da neutralidade de rede;

% V - preservação da estabilidade, segurança e funcionalidade da rede, por meio de medidas técnicas compatíveis com os padrões internacionais e pelo estímulo ao uso de boas práticas;

% VI - responsabilização dos agentes de acordo com suas atividades, nos termos da lei;

% VII - preservação da natureza participativa da rede;

% VIII - liberdade dos modelos de negócios promovidos na internet, desde que não conflitem com os demais princípios estabelecidos nesta Lei.
% \end{quoting}

\begin{quoting}
% Art. 11. Em qualquer operação de coleta, armazenamento, guarda e tratamento de registros, de dados pessoais ou de comunicações por provedores de conexão e de aplicações de internet em que pelo menos um desses atos ocorra em território nacional, deverão ser obrigatoriamente respeitados a legislação brasileira e os direitos à privacidade, à proteção dos dados pessoais e ao sigilo das comunicações privadas e dos registros.

% § 1º O disposto no caput aplica-se aos dados coletados em território nacional e ao conteúdo das comunicações, desde que pelo menos um dos terminais esteja localizado no Brasil.

% § 2º O disposto no caput aplica-se mesmo que as atividades sejam realizadas por pessoa jurídica sediada no exterior, desde que oferte serviço ao público brasileiro ou pelo menos uma integrante do mesmo grupo econômico possua estabelecimento no Brasil.
Art. 7º O acesso à internet é essencial ao exercício da cidadania, e ao usuário são assegurados os seguintes direitos:

[...]

VII - não fornecimento a terceiros de seus dados pessoais, inclusive registros de conexão, e de acesso a aplicações de internet, salvo mediante consentimento livre, expresso e informado ou nas hipóteses previstas em lei; 
\\
VIII - informações claras e completas sobre coleta, uso, armazenamento, tratamento e proteção de seus dados pessoais, que somente poderão ser utilizados para finalidades que:

a) justifiquem sua coleta;

b) não sejam vedadas pela legislação; e

c) estejam especificadas nos contratos de prestação de serviços ou em termos de uso de aplicações de internet; 
\end{quoting}

Esses dois incisos colocam algumas restrições sobre a coleta de dados pessoais.
Segundo ele, e o resto da legislação, não sendo explicitamente vedada por lei a circunstância da coleta, ela pode ser realizada contanto haja consentimento do usuário e transparência quanto ao que está sendo coletado.
Assim, respeitadas as boas-práticas de transparência e segurança desses dados delimitadas na lei, a coleta de dados pode ocorrer livremente desde que seja voluntária.

A LGPD segue uma linha muito parecida, exigindo que a coleta ocorra com pleno consentimento do usuário, com justificativa e publicidade do fato, e somente quando ela é estritamente necessária.
A maioria das restrições da legislação concentram-se no tratamento de dados sensíveis, como nome, endereço, documento de identidade, foto, localização, histórico de saúde, entre outros, mas, desde que o ente coletor se adeque as exigências postas acima, seja juridicamente imputável pelas suas ações, e mantenha um padrão de segurança suficiente, não existe proibição de coleta.
E nem seria isso uma estratégia viável para restringir o poder das Big Techs, uma vez que a maioria dos dados colhidos são registros banais, a partir dos quais, como foi dito no segundo capítulo, dados sensíveis podem ser reconstituídos.

O maior problema no que diz respeito a regulamentação do espaço digital é a total falta de transparência com que ele é conduzido.
Embora haja alguns impedimentos técnicos para ela poder se efetuar completamente, como os modelos ``caixa-preta'' discutidos no capítulo anterior, é necessário, dado que a coleta de dados é em parte necessária para sustentar os serviços digitais hoje funcionantes, que se tenha maior publicidade quanto a \textit{como} está ocorrendo o processamento desses dados, e, principalmente, \textit{para que fins}.
As últimas propostas legislativas que foram nessa linha, como o PL 2630/2020 (o ``PL das \textit{Fake News}''), foram duramente rejeitadas pelas grandes empresas de tecnologia.

Uma das peculiaridades da legislação brasileira é a territorialidade dela.
Segundo o décimo primeiro artigo do MCI, em ``qualquer operação de coleta, armazenamento, guarda e tratamento de registros, de dados pessoais ou de comunicações por provedores de conexão e de aplicações de internet em que pelo menos um desses atos ocorra em território nacional, deverão ser obrigatoriamente respeitados a legislação brasileira e os direitos à privacidade, à proteção dos dados pessoais e ao sigilo das comunicações privadas e dos registros.''

Tradicionalmente, o Estado é a autoridade máxima \textit{em seu território}, isso é, tem soberania interna.
Essa soberania interna é garantida e legitimada por uma soberania externa, um sistema internacional onde cada Estado reconhece a soberania interna de seus pares, abstendo-se de qualquer tentativa de sobrepôr-se a ele dentro de seu território.
A existência soberana de um Estado depende, portanto, de um território físico bem delimitado.
Qualquer disputa \textit{entre} Estados -- em outras palavras, disputas internacionais -- vive, nesse sistema, em anarquia ou anomia, completamente desregrada, e sem que nenhuma autoridade supranacional possa intervir, uma vez que não há nada acima do Estado. \cite{araujo_estado_2021}

A contradição provocada pela Internet é que, subitamente, os Estados precisam regular algo cuja existência é majoritariamente intangível (e, portanto, não existe fisicamente em seus territórios) e de natureza internacional.
Pela Internet alocar-se nessa área indefinida entre Estados, somada ao gigantesco poder internacional formado ao redor dela na forma das Big Techs e a dificuldade técnica em mediar o volume gigantesco de transações que por ela acontecem, é que a regulamentação da Internet é um tema tão difícil.
Isso abriu a possibilidade dessas empresas multinacionais, cuja origem confunde-se com a história de origem da Internet, ocuparem um importante papel socioeconômico em múltiplos países simultaneamente, com uma ``cadeia produtiva'' muito mais simplificada do que outros empreendimentos que requerem transporte físico entre territórios.

Hoje, é quase impossível imaginar o Brasil funcionando sem acesso ao WhatsApp, por exemplo.
Isso significa que a sociedade brasileira vive em uma relação de dependência internacional em um âmbito econômico, social e tecnológico.
A abertura da Internet no Brasil atropelou qualquer iniciativa que pudesse existir aqui com empresas estrangeiras muito mais consolidadas, e que tiveram uma vantagem inicial ao virem de países mais industrializados.
Muito se fala do Estado chinês como um exemplo de censura e fechamento do espaço virtual para o exterior -- o tal ``Grande Firewall da China''.
Porém, para além de uma medida de censura essa também é uma medida de protecionismo econômico, permitindo o surgimento de uma indústria local muito mais forte do que no caso brasileiro, e impedindo a consolidação de empresas estrangeiras como uma potente força de \textit{lobby}.
